Human-grounded validation remains a central gap in large language model (LLM) evaluation for affective text understanding. Existing benchmarks and multilingual profiling can reveal model-to-model or language-to-language variation, but such variation alone does not show which systems are closest to human judgments. This paper proposes a compact human-grounded multilingual benchmark for literary emotion evaluation using three Japanese literary works and four continuous affect dimensions: interest, surprise, sadness, and anger. Japanese Visual Analogue Scale (VAS) ratings from student readers define the reference profile, and six contemporary LLMs are compared through a unified multi-provider experimental platform. The primary endpoint is alignment with the Japanese human reference under a neutral-reader prompt; English and Chinese conditions are treated as robustness tests for cross-language drift, not as co-equal human-grounded endpoints. Claude Sonnet 4.5 achieved the lowest mean absolute error to the Japanese human reference (20.026), whereas Gemini 2.5 Pro showed the largest error (28.952). GPT-5.4 achieved the strongest rank-order agreement with the human profile, while Qwen3.6 Plus and GPT-5.4 exhibited the smallest average cross-language drift (3.346 and 3.354, respectively). The results show that primary human alignment and multilingual robustness are related but separable properties, making the benchmark suitable for compact and reproducible model screening.
